\documentclass[a4paper,11pt,fleqn]{article}%fleqn pour aligner les equations à gauche
\usepackage{../../Styles/Cours}

\newcommand{\chnum}{Chapitre 2}
\newcommand{\chtitle}{Échelle de teinte}
\newcommand{\dtype}{TP}


\begin{document}
\maketitle

Dans une course un cycliste vient de faire une chute. Sa jambe a frotté contre le bitume et présente une plaie importante avec un saignement. Le soigneur arrive sur les lieux et le prend en charge.\\
\vspace{.5cm}

\noindent
\fbox{
\begin{minipage}{.98\textwidth}
\textbf{Document 1 : Le Dakin}\\
La ligueur de Dakin (ou eau de Dakin) est un liquide antiseptique (bactéricide, fongicide et virucide) utilisé pour le lavage des plaies et de muqueuses, de couleur rose et à l’odeur d’eau de Javel. Son action se base sur l’hypochlorite de sodium (eau de Javel diluée) additionnée de permanganate de potassium pour le stabiliser vis-à-vis de la lumière. C’est le permanganate de potassium qui donne à l’eau de Dakin sa coloration rose. La solution doit être conservée à l’abri de la lumière pour ralentir sa décomposition qui est rapide (quelques jours).\\
Un préparateur en pharmacie a réalisé une solution $S_0$ de permanganate de potassium de concentration en masse $\gamma$ = \num{0,03} \unit{\gram \per \liter}.
\end{minipage}
}

\vspace{.5cm}
\noindent
\fbox{
\begin{minipage}{.98\textwidth}
\textbf{Document 2 : Principe de la manipulation}\\
On peut évaluer la concentration en masse d’une espèce chimique colorée dans une solution grâce à la couleur de cette solution. A partir d’une solution $S_0$ contenant une espèce chimique colorée, on prépare des solutions diluées de différentes concentrations. On réalise ainsi une « échelle de teinte » qui va servir, par comparaison de la couleur, à évaluer la concentration inconnue d’une solution $S$.
On cherche ici à déterminer la concentration de la solution de Dakin, une solution antiseptique de couleur violette contenant du permanganate de potassium.
\textbf{Rappel :} la solution $S_0$ est appelée « solution mère » et les solutions diluées sont appelées « solutions filles »
\end{minipage}
}
\vspace{.5cm}

\noindent
\fbox{
\begin{minipage}{.98\textwidth}
\textbf{Document 3 : Protocole expérimental}
\begin{itemize*}
	\item Verser une quantité de solution $S_0$ dans un bécher
	\item Utiliser une pipette (jaugée ou graduée) équipée d’une propipette pour prélever le volume demandé de solution $S_0$.
	\item Transférer le volume prélevé dans une fiole jaugée de 50,0 mL
	\item Compléter la fiole jaugée jusqu’au trait de jauge
	\item Agiter
\end{itemize*}
\end{minipage}
}
\vspace{.5cm}

\noindent
\fbox{
\begin{minipage}{.98\textwidth}
\textbf{Document 4 : Calcul du facteur de dilution}\\
Le facteur de dilution d’une solution noté $F$, est le rapport entre la concentration de la solution mère $\gamma_0$ et de la solution fille $\gamma_i$. Il est égal à $F = \dfrac{\gamma_0}{\gamma_i}$ ou encore $F = \dfrac{V_i}{V_O}$
\end{minipage}

}
\newpage

\noindent
\fbox{
\begin{minipage}{.98\textwidth}
\textbf{Document 5 : Solutions à préparer}\\
\begin{tabular}{{|>{\centering}m{3.5cm}|>{\centering}m{2cm}|>{\centering}m{2cm}|>{\centering}m{2cm}|>{\centering}m{2cm}|>{\centering}m{2cm}|>{\centering\arraybackslash}m{2cm}|}}
\hline
\textbf{Solution}&$S_0$&$S_1$&$S_2$&$S_3$&$S_4$&$S_5$\\
\hline Volume $V_0$ de solution $S_0$ prélevé (\unit{\milli \liter}) &\num{50,0}&\num{25,0}&\num{20,0}&\num{12,5}&\num{10,0}&\num{5,0}\\
\hline Volume final de la colution $V_i$&\num{50,0}&\num{50,0}&\num{50,0}&\num{50,0}&\num{50,0}&\num{50,0}\\
\hline Facteur de dilution $F$& & & & & & \rule[1cm]{0cm}{0cm}\\
\hline Concentration en masse de soluté $\gamma$ \unit{\gram \per \liter}& & & & & & \\
\hline
\end{tabular}
\end{minipage}
}
\begin{enumerate}
	\item Compléter le tableau pour toutes les solutions en détaillant les calculs dans chaque cas.
	\item Comparer la solution inconnue aux solutions de l’échelle de teinte et en déduire un encadrement de sa concentration.
	\item L’étiquette du flacon de Dakin mentionne « pour 100 mL de solution, permanganate de potassium : \num{0,010} \unit{\gram} ». Calculer la concentration en masse de permanganate de potassium de la solution de Dakin.
	\item La valeur calculée est elle en accord avec celle de la concentration estimée à l’aide de l’échelle de teinte ?
\end{enumerate}

\end{document}